\section{The model}
\label{sec:model}

Everything in this report is stated over the objects defined here. Each
definition transcribes a numbered section of \SPEC{} and cites it. Where
the transcription makes a choice the source leaves implicit, the choice is
flagged in a remark rather than buried.

\subsection{History}

\begin{definition}[commit history, \spec{12.2}]
\label{def:history}
The history is a finite directed acyclic graph on a set of commits \Cmts, with
an edge from a commit to each of its parents. A distinguished commit
$\Head \in \Cmts$ is \emph{HEAD}. For $c \in \Cmts$,
$\Rch{c} \subseteq \Cmts$ is the set of commits reachable from $c$ by following
parent edges, including $c$ itself. Write $c' \anc c$ for $c' \in \Rch{c}$
(``$c'$ is an ancestor-or-self of $c$'').
\end{definition}

\Head{} is fixed for the whole report. The retry and convergence guarantees of \spec{13.7c} are stated at a fixed \Head{}; advancing it may legitimately add, remove, or redirect work. Structural termination and single-state determinism do not need that restriction.

\begin{remark}
Ancestry, not commit date or log order, is the only ordering used. This is what
makes cancellation deterministic under merges and rebases (\spec{10.4}) and
what lets the model ignore clocks entirely.
\end{remark}

\subsection{Workspace}

\begin{definition}[workspace graph, \spec{13.1}, \spec{8.4}]
\label{def:graph}
\Pkgs{} is the finite set of packages present at \Head{}. For a fixed set of
edge kinds $K$ (the configured \code{propagation.kinds}), $\Edges \subseteq
\Pkgs \times \Pkgs$ contains $(a,b)$ when $a$ declares a dependency on $b$ of a
kind in $K$. The \emph{dependents} of $p$ are $\Dep{p} = \{d : (d,p) \in
\Edges\}$. The graph $(\Pkgs, \Edges)$ is acyclic.
\end{definition}

Acyclicity is an assumption the model may make because the engine enforces it:
a cycle over the propagation or publish-ordering kinds is the repository-scoped
error \diag{E200} and aborts the run before any plan is computed. \spec{13.1}
justifies rejecting rather than accommodating cycles: no publish order exists
for a cycle, so whichever member goes first declares a range on a version that
does not yet exist, and if the run then fails that state is permanent.

Propagation travels \emph{against} the dependency edges, from a changed package
to the packages that consume it; that is, along \Dep{\cdot}.

\begin{example}[the running workspace, \spec{22}]
\label{ex:workspace}
The running examples use the six-package workspace shared by many conformance vectors:
\[
  \Pkgs = \{\,\code{core},\ \code{cli},\ \code{ui},\ \code{api},\
            \code{docs-site},\ \code{@acme/theme}\,\},
\]
\[
  \Edges = \{\,(\code{cli},\code{core}),\ (\code{ui},\code{core}),\
              (\code{api},\code{core}),\ (\code{docs-site},\code{ui}),\
              (\code{@acme/theme},\code{ui})\,\},
\]
drawn in \cref{fig:workspace}, with the initial tag state
\[
  \st_0 = \{\ \code{core@1.4.2},\ \code{cli@2.0.0},\ \code{ui@0.9.1},\
              \code{api@1.2.0},\ \code{@acme/theme@1.0.0}\ \}.
\]
\code{docs-site} has never been released, so it has no baseline; it is private
and publishes to an internal registry, which by \spec{13.10a} exempts it from
nothing.
\end{example}

\begin{figure}[t]
\centering
\begin{tikzpicture}[
    >={Stealth[width=5pt,length=5pt]},
    pkg/.style={draw,rounded corners=2pt,inner sep=4pt,font=\small\ttfamily},
    every edge/.style={draw,->}]
  \node[pkg] (core) at (0,0) {core};
  \node[pkg] (cli)  at (-2.6,1.5) {cli};
  \node[pkg] (ui)   at (0,1.5)    {ui};
  \node[pkg] (api)  at (2.6,1.5)  {api};
  \node[pkg] (theme) at (-1.5,3.0) {@acme/theme};
  \node[pkg] (docs)  at (1.5,3.0)  {docs-site};
  \draw[->] (cli)   -- (core);
  \draw[->] (ui)    -- (core);
  \draw[->] (api)   -- (core);
  \draw[->] (theme) -- (ui);
  \draw[->] (docs)  -- (ui);
  \node[font=\footnotesize,anchor=west] at (3.2,0.75)
    {$A \rightarrow B$: $A$ depends on $B$};
  \node[font=\footnotesize,anchor=west] at (3.2,0.15)
    {propagation runs upward};
\end{tikzpicture}
\caption{The running workspace of \cref{ex:workspace}. Distances from
$S=\{\code{core}\}$ are $\dist{S}{\code{cli}} = \dist{S}{\code{ui}} =
\dist{S}{\code{api}} = 1$ and $\dist{S}{\code{docs-site}} =
\dist{S}{\code{@acme/theme}} = 2$.}
\label{fig:workspace}
\end{figure}

\begin{definition}[distance, \spec{9.2}]
\label{def:dist}
For $S \subseteq \Pkgs$ and $d \in \Pkgs$, $\dist{S}{d}$ is the least $n$ for
which there are $p_0 \in S$ and $p_1,\dots,p_n = d$ with $p_{i+1} \in
\Dep{p_i}$ for all $i$, and $\infty$ if no such chain exists. Note
$\dist{S}{d} = 0$ exactly when $d \in S$.
\end{definition}

\subsection{Bumps and versions}

\begin{definition}[bumps, \spec{7.1}, \spec{9.1}]
\label{def:bump}
$\Bumps = \{\bnone, \bpatch, \bminor, \bmajor\}$ ordered
$\bnone < \bpatch < \bminor < \bmajor$. Being a finite chain, $(\Bumps, \leq)$
is a join-semilattice with least element \bnone{} and join $\sqcup = \max$.
\end{definition}

The whole of \spec{9.1}'s ``bumps accumulate'' is the statement that the
combining operator is this join. Two consequences that \SPEC{} states
separately are then immediate and are used without further comment: combination
is associative, commutative and idempotent, so the order in which contributions
are gathered cannot matter; and $x \sqcup y \geq x$, so propagation can only
ever raise a package's bump, never lower or replace a direct decision.

\begin{definition}[versions and channels, \spec{11.1}--\spec{11.4}]
\label{def:version}
\Vers{} is the set of SemVer 2.0.0 versions with SemVer precedence. $\Chans =
\{\stable\} \cup \mathcal{N}$ for a set $\mathcal{N}$ of channel names.
$\chanof{v} = \stable$ if $v$ carries no prerelease identifier, and otherwise
the first (alphanumeric) prerelease identifier of $v$. $\coreof{v}$ is the
\texttt{major.minor.patch} triple of $v$ and $\cntof{v}$ the numeric counter of
a prerelease. $\applyb{v}{b}$ applies bump $b$ to $v$, with $\applyb{v}{\bnone}
= v$.
\end{definition}

\subsection{State}

Everything above is fixed by \Head{} and the configuration. Exactly one thing
changes from run to run, and isolating it is what makes the guarantees provable
as statements about a sequence.

\begin{definition}[engine state, \spec{13.2}, \spec{12.3}]
\label{def:state}
A state \st{} is a finite set of release tags $(p, v, c) \in \Pkgs \times \Vers
\times \Cmts$ with $c \anc \Head$. States are ordered by inclusion,
$\st \leq \st'$ iff $\st \subseteq \st'$. From \st{}:
\begin{itemize}
  \item $\base{\st}{p}$ is the highest-precedence version among the tags of $p$
    in \st{}, undefined if there are none;
  \item $\basec{\st}{p}$ is the commit of that tag;
  \item $\stbase{\st}{p}$ is the highest-precedence non-prerelease version among
    them, and $\stcommit{\st}{p}$ its commit.
\end{itemize}
Where $\base{\st}{p}$ is undefined, $\chanof{\base{\st}{p}} = \stable$ by
convention (\spec{11.1}): a never-released package is graduated by nothing and
is put on a train by any directive naming it.
\end{definition}

The baseline functions additionally assume a unique highest-precedence
reachable tag for each package and each stable baseline. States with tied
precedence and different commits need a specified tie rule and are outside
this proof's domain. Operational lock and temporary refs are not release tags
and are not members of $\st$.

That release-tag states only ever grow is not an assumption about tidy operators; it is
normative. \spec{19.1} forbids moving or deleting a pushed tag and \spec{12.4}
makes tags authoritative over manifests, because a tag is the record that a leg
of the transaction completed and an unwritten-back or rewritten record is
exactly what makes a saga unrecoverable.

\begin{definition}[train and undelivered windows, \spec{13.3}]
\label{def:window}
\begin{equation}
  \Win{\st}{p} \;=\; \Rch{\Head} \setminus \Rch{\stcommit{\st}{p}},
  \label{eq:window}
\end{equation}
and $\Win{\st}{p} = \Rch{\Head}$ when $p$ has no stable baseline in \st{}.
The undelivered window is
\[
  \Winf{\st}{p}=\Win{\st}{p}\setminus\Rch{\basec{\st}{p}},
\]
with an empty subtracted set when $p$ has no baseline tag.
\end{definition}

$\Win{\st}{p}$ is the train window used to aggregate a prerelease target from
its stable base. $\Winf{\st}{p}$ is the fresh-admission window for direct
work, corrections, channels, and propagation to dependent $p$. For a stable
baseline they coincide; on a train they must not be conflated.

\begin{example}[windows]
\label{ex:window}
Let the history be $c_0 \to c_1 \to c_2 = \Head$ with $\st_0$ of
\cref{ex:workspace}, where the tags sit at $c_0$. Then
\[
  \Win{\st_0}{\code{core}} = \{c_1, c_2\},
  \qquad
  \Win{\st_0}{\code{docs-site}} = \{c_0, c_1, c_2\},
\]
the second because a package with no stable baseline has the whole reachable
history pending: \spec{12.5} makes an unreleased package's first release
account for everything, which is also why adopting the engine on an existing
repository does not need a migration step.
\end{example}

\begin{remark}[the window is measured from the last \emph{stable} tag]
\label{rem:stable-window}
This is deliberate and it serves \spec{11.4}: a package on a prerelease train
must accumulate every commit since its last \emph{stable} release, because the
train's target version is recomputed from the stable baseline on every run so
that a breaking change arriving mid-train moves the whole train. The
consequence, which \cref{sec:obligations} returns to at length, is that a
commit already published in a prerelease of $p$ is still a member of
$\Win{\st}{p}$.
\end{remark}

\subsection{Units}

\begin{definition}[units, \spec{4.4}, \spec{13.4}, \spec{13.4b}]
\label{def:unit}
\Units{} is the set of units carried by commits in $\bigcup_{p}\Win{\st}{p}$,
after parsing (\spec{20}), scope resolution (\spec{6}) and the application of
corrections (\spec{13.4b}). For $u \in \Units$:
$\cof{u} \in \Cmts$ is its commit; $\scopef{u} \subseteq \Pkgs$ its resolved
scope-set; $\bumpof{u} \in \Bumps$ the bump its type and \code{!} map to; and
its directives are
\[
  \pbump{u} \in \Bumps,\quad
  \pdepth{u} \in \mathbb{N} \cup \{\infty\},\quad
  \pscope{u} \subseteq \Pkgs,\quad
  \pchan{u},\quad
  \pcdepth{u},\quad
  \pcscope{u} \subseteq \Pkgs,
\]
where $\pbump{u}$ is the \emph{effective} propagated bump: the value the unit
states, or $\bumpof{u}$ where it states \code{inherit}. The two axes default
differently and the asymmetry is deliberate (\spec{8.1}). \code{Propagate}
defaults to \bpatch, not to \code{inherit}: a dependent reached by propagation
has had a dependency change under it but no change to its own public API, so
\bpatch{} is the honest signal, and \code{feat(core)\^{}} therefore gives its
consumers a patch rather than a minor. \code{Propagate-Channel} defaults to
\code{inherit}, since a channel is not a quantity that can be understated.
Both defaults apply only once the corresponding depth is non-zero, and both
depths default to $0$ (\spec{8.3}, \spec{8.3a}): a unit that says nothing
releases its own packages and nothing else. $\pchan{u}$ is a channel value: a
name, \stable, \code{inherit}, \code{none}, or a transition
$\langle f \to t\rangle$.
\end{definition}

Taking corrections as already applied is what \spec{13.4b} prescribes: it
requires \code{Edits} and \code{Deletes} to be resolved into the tuple stream
before cancellation runs, and notes that every later phase operates on the
corrected stream, which downstream phases cannot distinguish from one that was
authored that way.

\begin{remark}[the corrected stream depends on \st{}]
\label{rem:corrections}
It does not follow, and \spec{13.4b} does not claim, that the stream is
independent of the state. A correction reaches only \emph{undischarged} work:
it applies to a $(\text{package}, \text{unit})$ pair only while that package's
window still contains the target commit, and once the package has released the
target the correction is a no-op for it (\diag{W209}, \spec{7.4.2}). So
$\Units{}$ is properly a function $\Units_{\st}$, and a record suppressed by a
\code{Deletes} for a package returns to the stream for that package once the
package releases. This matters in exactly one place, and it is the same place
that a lapsed cancellation matters: \cref{ob:g5}. Everywhere else the two
readings coincide, because a released package's own window is empty and its
direct contribution is therefore already spent, so the notation
$\Units{}$ is kept unindexed below.
\end{remark}

\begin{definition}[suppression and source packages, \spec{13.4a}]
\label{def:sources}
$\canc{c}{p}$ holds when $c$ is an ancestor-or-self of some \code{cancel} commit
whose resolved scope contains $p$ (\spec{10.3}); $\held{p}$ holds when $p$'s
effective \code{Release-As} directive is \code{none} (\spec{13.6a}). Both are
computed from the corrected stream, \Head{}, configuration, and release-tag state. They are constant across retries only where the retry invariant says those inputs remain fixed. Write
\[
  \disch{\st}{p}{c} \;\equiv\; c \notin \Winf{\st}{p},
\]
and then
\[
  \srcp{\st}{u} \;=\;
    \bigl\{\, p \in \scopef{u} \;:\;
       \disch{\st}{p}{\cof{u}} \;\lor\;
       \lnot\bigl(\canc{\cof{u}}{p} \lor \held{p}\bigr) \,\bigr\}.
\]
\end{definition}

The disjunction is the normative content of \spec{13.4a}: suppression applies
only to \emph{undischarged} work. Once $p$ has published the version carrying
$u$, the artefact its consumers are owed is public, and neither a later
\code{cancel} (a no-op, \diag{W170}) nor a later hold may retroactively strand
those consumers. \Cref{ob:g5} shows that this clause, and not the traversal, is
where \G{5} is decided.

\subsection{Runs}

\begin{definition}[plan]
\label{def:plan}
Given \st{}, the release computation of \spec{13} yields the bumps
$\dirb{\st}{p}$ and $\propb{\st}{p}$ (\cref{def:direct,def:propagated}), the
channels $\chanres{\st}{p}$ (\cref{def:channels}), and
$\effb{\st}{p} = \dirb{\st}{p} \sqcup \propb{\st}{p}$. The plan is
\begin{align*}
  \plan{\st} \;=\; \bigl\{\, p \in \Pkgs \;:\;
    &\lnot\held{p} \;\land\;
    \bigl(\effb{\st}{p} \neq \bnone \\[-1mm]
    &\;\lor\; \chanres{\st}{p} \neq \chanof{\base{\st}{p}}
    \;\lor\; p \text{ has an applicable exact } \code{Release-As} \\[-1mm]
    &\hspace{35mm}\text{in }\Winf{\st}{p} \bigr) \,\bigr\},
\end{align*}
each $p \in \plan{\st}$ carrying the version $\vers{\st}{p}$ of \spec{13.9}. The fresh bumps above decide whether work is newly admissible; they are not by themselves the prerelease target. On a train, the intended calculation aggregates the surviving corrected units in $\Win{\st}{p}$ from the stable baseline, including work already delivered to an earlier prerelease. For graduation, the model applies the specific rule of \spec{11.5}: publish that same train target without a prerelease suffix. The generic summary in \spec{13.9} is inconsistent with this rule; \cref{ob:graduation} records the textual obligation.
\end{definition}

\begin{remark}[\plan{\st} is the released set, not the printed one]
\label{rem:plan-held}
\spec{13.9} computes a would-be version for a held package too, and \spec{13.6a}
requires the engine to report it (\diag{W154}) so that the number needed to lift
the hold is available without hand computation. That reported entry is part of
what an operator calls ``the plan'', but it is not a release: \spec{13.10}
excludes it. \plan{\st} above is the released set, which is why \cref{thm:g6}
needs no exception for holds and why the convergence check of \spec{19.6} states
one: the two are reading the same word differently, and only \spec{19.6}'s
reading has to be told about held packages.
\end{remark}

\begin{definition}[run, \spec{19.1}--\spec{19.3}]
\label{def:run}
A run from \st{} is parameterised by a set $F \subseteq \plan{\st}$ of
publishes that fail. It induces the blocked set
\begin{align*}
  B \;=\; \{\, p \in \plan{\st}\setminus F \;:\;
    &\ p \text{ transitively depends on some } q \in F \\[-1mm]
    &\ \text{over } \code{publish.blockingKinds} \,\},
\end{align*}
publishes $\plan{\st} \setminus (F \cup B)$ in the order of \spec{19.2}, and
yields
\[
  \Run{\st}{F} \;=\; \st \;\cup\;
  \bigl\{\,(p, \vers{\st}{p}, \Head) \;:\; p \in \plan{\st}\setminus(F\cup B)\,\bigr\}.
\]
\end{definition}

Three features of this definition carry weight later.

$F$ is chosen adversarially and independently on each run. The guarantees are
claims about \emph{every} failure schedule, not about a well-behaved one; a
proof that appeals to what usually fails proves nothing about a system whose
purpose is to survive what does not usually happen.

Every tag is written at \Head{}, never at the commit the work arrived in
(\spec{19.1}). This is what makes \cref{lem:window} a two-case statement rather
than an argument about which commits a tag covers.

A package is tagged if it is released, whatever its publish target, including
targets of \code{none} that upload nothing (\spec{13.10a}). \Cref{prop:tagall}
shows this is forced: the alternative is a permanent hole in \G{6}.

\begin{definition}[run sequence]
\label{def:seq}
A \emph{run sequence} at \Head{} is $\st_1 \leq \st_2 \leq \cdots$ with
$\st_{k+1} = \Run{\st_k}{F_k}$ for some $F_k \subseteq \plan{\st_k}$. By
\cref{def:run} it is monotone: $\st_k \subseteq \st_{k+1}$.
\end{definition}
